% -----------------------------------------------------------------------------
% Section 2: Related Work
% Paper: "Silent Decisions Are Type Errors: Enforcing AI Accountability
%         via Linear Resource Types"
% Venue: IEEE Computer — Special Issue on AI Governance
% Deadline: 13 April 2026
% Authors: Daniel Lau Pereira Soares
% Rev 3: 25 March 2026 — ACID distinction added (§2.2)
% -----------------------------------------------------------------------------

\section{Related Work}
\label{sec:related_work}

Algorithmic accountability intersects infrastructure governance, runtime
verification, and type theory.  We situate our contribution against the
prevailing methodology in each domain.

% ----------------------------------------------------------------
\subsection{Policy-as-Code for Infrastructure Governance}
\label{sec:pac}

The industry standard for continuous compliance relies on
Policy-as-Code (PaC) frameworks such as Open Policy
Agent~\cite{opa2024}, Kyverno, and Checkov.
These tools excel at governing infrastructure artifacts---that a
Kubernetes pod lacks privilege escalation, or that a Terraform plan
avoids over-privileged IAM roles.
However, they are designed to evaluate \emph{declarative
configurations}, not \emph{dynamic runtime decisions} affecting
individual rights.

Evidence of compliance manifests as asynchronous audit logs: records
written after the decision logic executes, stored separately, and
subject to omission or alteration if the logging pipeline fails.
PaC frameworks govern configuration artifacts, not transient decision
acts with legal consequences for a person; accountability remains a
post-hoc state, structurally vulnerable to the silent-decision
failure mode of Section~\ref{sec:intro}.

% ----------------------------------------------------------------
\subsection{Runtime Assertions and Audit Logging}
\label{sec:runtime}

To govern application-level decision logic, practitioners
typically rely on a combination of runtime assertions
(e.g., \texttt{assert(evidence != null)}) and structured
audit logging.
While functional under nominal conditions, these are procedural
promises, not structural invariants.
Assertions can be disabled in production builds (stripping
\texttt{debug\_assertions} in Rust; omitting \texttt{-ea} in the
Java VM); audit logs can be truncated under load, dropped by a
misconfigured sink, or bypassed by an untested code path.

Both mechanisms operate \emph{after} the decision logic executes:
they do not prevent a verdict that lacks integrity-authenticated
evidence from materializing; they attempt to record its absence, if
they run at all.
This architectural decoupling separates the act of deciding from the
act of preserving evidence needed to support review, explanation, and
oversight under LGPD Art.~20 and EU AI Act Arts.~14 and~86.

A stronger form of runtime guarantee is offered by transactional
databases: an ACID transaction wrapping both the decision record and
the audit log ensures that either both are committed or neither is.
However, transactional atomicity is a tool the developer must
consciously choose to apply---wrapping decision and log within the
same \texttt{BEGIN}/\texttt{COMMIT} boundary.
Omitting that boundary is a logic error, not a type error: the
compiler accepts the program, the database commits the decision
without the log, and the silent-decision failure mode materialises
with no static diagnostic.
BTV closes exactly this gap: the compiler rejects any program that
constructs a \texttt{Verdict} without a causally-prior
\texttt{EvidenceToken}, regardless of whether the developer intended
to include the evidence or not.

% ----------------------------------------------------------------
\subsection{Linear Types and Resource Correctness}
\label{sec:linear_related}

Linear Logic's structural resource guarantees~\cite{girard1987linear,wadler1991linear}
were adapted into programming-language design and underpin Rust's
ownership and borrowing model~\cite{jung2017rustbelt}.
Prior applications of linear and substructural types, however, focus
predominantly on \emph{resource} correctness: ensuring that memory is
freed exactly once, that file handles are not leaked, or that a
communication protocol is followed in order.  Vault enforces
high-level resource-management protocols in low-level code via
linear typestates~\cite{deline2001vault}; typestate systems make
state-dependent protocol obligations explicit at the type
level~\cite{pierce2002tapl}; session types impose
communication protocols on channels; and the substructural design
space of affine versus linear versus relevant systems is surveyed by
Walker~\cite{walker2005substructural}, with state-dependent
disciplines studied by Ahmed et al.~\cite{ahmed2009state}.
Our claim C1 is deliberately narrower: we \emph{instantiate} that
established discipline over \emph{regulatory-support} artifacts---a
BTV verdict must consume causally bound evidence and an
authority-issued compliance token---and relate the mechanism to LGPD
Art.~20 and EU AI Act Arts.~14 and~86.
We scope the resulting guarantee to the Safe Rust
type perimeter (Section~\ref{sec:encapsulation}), with the escape
hatches of Section~3 stated explicitly.
Bound-To-Verdict extends the linear-type discipline from resource
management to regulatory accountability, not from nothing to
regulatory accountability.
